\chapter{基因组会复制、借用和改造}

\begin{kaichang}
秋冬的街角，烤红薯的炉子又支起来了。掰开一个烤得流糖的红薯，咬一口，又甜又面。你可能不知道：这个红薯，严格说来是一个\textbf{天然的转基因生物}。

2015 年，一群科学家检测了来自世界各地的近三百个栽培红薯品种，结果在\textbf{每一个}品种的基因组里，都发现了同一种土壤细菌——农杆菌——留下来的基因片段。也就是说，在很久以前，农杆菌把自己的一段 DNA“塞”进了红薯祖先的细胞，这段外来基因不但没有惹祸，反而被红薯留了下来，一直传到今天每一块烤红薯里。

没有实验室，没有科学家，没有移液枪。大自然自己就干过“转基因”这件事。那么，基因组到底是怎样一台机器？它不是应该被严格保护、逐字照抄的生命蓝图吗？这一章，我们来看看这本“生命之书”真实的样子：它会被整页影印、会自己搬家、会向邻居借页，还会在旧零件上敲敲打打，改造出全新的东西。
\end{kaichang}

\section{一本并不安分的书}

到目前为止，这本书给你建立的图景可能是这样的：DNA 是一份精密的说明书（第 70 章），复制时会有抄写员校对（第 72 章），突变大多是拼写错误——换掉一个字母、漏掉一个字母。这些都没错，但它们只描述了“改一个字母”这种最细小的变化。

如果你只盯着单个字母，就会错过这本“书”更戏剧性的一面。真实的基因组里，时刻发生着大尺度的变动：

\begin{itemize}[itemsep=2pt]
\item 整整一段文字被\textbf{影印}成两份、三份、甚至几十份；
\item 某些段落\textbf{自己搬家}，从书的这一章跳到那一章；
\item 有时候干脆\textbf{从别的书里撕一页}，直接粘进来；
\item 还有亿万年前“死掉的病毒”留下的\textbf{残页}，夹在书里，有的竟然被重新启用。
\end{itemize}

这听起来像图书管理员的噩梦。但请记住第 76 章讲过的：进化的原料就是变异。字母拼错是一种变异，而影印、搬家、借页，是规模更大的变异。这一章的核心问题是：这些“大动作”怎样发生？它们大多带来什么后果？又有哪一些，碰巧成了进化史上的关键发明？

先打个预防针：这些变化\textbf{不是}基因组“为了”进化而做的。没有任何分子知道未来需要什么。变化只是按化学和概率的规则不停地发生，然后由环境和时间来筛选——这条因果链的顺序，从头到尾都不能倒过来。

\section{影印一整段：基因重复}

先从一个你能亲眼验证（至少能对着镜子验证）的事实说起：你能分辨红色和绿色。而你家的狗，很可能分不清。

狗的世界大致是两种颜色调出来的，而你有三种“颜色传感器”——视网膜里三种视蛋白，分别对蓝、绿、红最敏感。决定红、绿视蛋白的两个基因，在 X 染色体上肩并肩躺着，序列相似度高达约 98\%。为什么这么像？因为它们本来是\textbf{同一个基因的两份拷贝}。

大约三四千万年前，我们灵长类的祖先身上发生了一次“影印事故”：在制造卵细胞或精子时，两条染色体配对、交换片段（第 69 章讲的减数分裂），偶尔会对不齐——一条染色体上的基因 A，对上了另一条上的基因 B，交换的结果，一个细胞多拿了一段，另一个少了一段。这种“错位交换”会把一整段 DNA 影印成两份。红绿视蛋白基因，就是一次古老影印留下的双胞胎。两份拷贝后来各自积累小突变，一个变得对红光敏感，一个保持对绿光敏感——我们祖先的世界里，从此多出了一个颜色维度。在只有两种视蛋白的哺乳动物背景里，能区分成熟红果和绿叶，是实打实的生存优势。

这就是\textbf{基因重复}（gene duplication）：一整段 DNA——有时是一个基因，有时是好几个基因——被影印出额外的拷贝。它是进化最重要的“原材料生产线”，道理很朴素：原件只有一份时，它忙着干自己的活，突变稍微大一点就可能把活搞砸，自然选择会立刻把这类突变清走；可一旦有了两份，一份继续干原来的活，另一份就“自由”了——它可以随便积累突变，绝大多数时候变成一个没用的残次品（假基因），但偶尔，它会撞出一个有用的新功能。

你身体里有一整个家族都是这样“影印再分化”出来的。第 54 章提到过的血红蛋白，其实有好几个版本：胎儿的血红蛋白和成人的不一样，前者从母亲的血液里“抢”氧气的本领更强。控制它们的 $\alpha$ 珠蛋白和 $\beta$ 珠蛋白基因家族，就是几十亿年间一轮轮影印、再各自改装的产物——肌肉里存氧气的肌红蛋白，是这条影印链上更古老的一位成员。

基因重复甚至会因为你吃的是什么而留下痕迹。唾液里的淀粉酶——就是第 56 章那个把淀粉嚼出甜味的酶——由 \emph{AMY1} 基因制造。不同人携带的 \emph{AMY1} 拷贝数差异很大：有人只有 2 份，有人多达十几份。研究发现，祖上长期食用大量淀粉（稻米、麦、薯类）的人群，平均拷贝数更高，口水里的淀粉酶也更多。换句话说，连“你碗里主要装什么”这种文化史，都在基因组上留下了影印次数的记录。

\begin{qiaoliang}
你的基因组里，到底有多少东西是“影印品”和“外来户”？拿人类单倍体基因组约 $3\times 10^{9}$ 个碱基对来算一笔账：

\begin{itemize}[itemsep=2pt]
\item 真正编码蛋白质的序列，只占约 $1.5\%$——大约 4500 万个碱基；
\item 由“会搬家的段落”（转座元件，下一节讲）构成的序列，占约 $45\%$——超过 13 亿个碱基；
\item 古代病毒留下的残页，占约 $8\%$——比全部编码蛋白质的部分还多五倍。
\end{itemize}

再算一个具体的：人类基因组里有一种叫 \emph{Alu} 的搬家元件，长度约 300 个碱基，拷贝数超过 100 万份。$10^{6}\times 300=3\times 10^{8}$，即 3 亿个碱基——光这一种元件，就占了整个基因组的约十分之一。所以“基因组是一本由基因组成的书”这个说法，离真实图景很远：它更像一本被反复影印、贴满便签、夹着大量旧广告单的工作手册，说明书正文只占其中薄薄几页。
\end{qiaoliang}

\section{会自己搬家的段落：转座元件}

影印一段 DNA，至少需要染色体错位交换这样的“大事故”。但基因组里还有一类更不安分的居民：它们自带搬家工具，能从书的一个位置剪下来（或者复制一份）贴到另一个位置。它们叫\textbf{转座元件}，俗名“跳跃基因”。

\begin{zhengju}
1940 年代，美国科学家芭芭拉·麦克林托克（Barbara McClintock）在冷泉港的试验田里种玉米。她盯上了一类奇怪的玉米粒：同一粒玉米上，紫色和白色的斑块犬牙交错，而且斑块的分布模式居然能遗传、能预测。

当时的标准解释是：色素基因不稳定，爱突变。但麦克林托克用显微镜观察玉米细胞的染色体，看到了别人忽略的细节：染色体上有一个特定位点总是发生断裂，而断裂的位置会随着世代\textbf{移动}。她提出一个离经叛道的假说：基因组里存在可以移动的“控制元件”（她命名为 \emph{Ds} 和 \emph{Ac}），它们跳到色素基因中间，色素就生产不出来，玉米粒出现白斑；跳走之后，细胞又恢复制造紫色。斑块，就是“跳跃”在细胞后代里留下的足迹。

这个假说在 1950 年代几乎没人信——教科书说基因组是稳定排列的珠子串，哪来的“跳”？麦克林托克一度懒得再发表。直到 1960 到 70 年代，分子生物学家在细菌里也找到了会移动的 DNA 片段，她的玉米实验才被重新翻出来。1983 年，81 岁的她独自获得诺贝尔奖。这个故事里最值得记住的是她的方法：她没有因为现象“怪”就扔掉它，而是把斑块的分布、染色体的断裂、世代的传递三条独立的线索对在一起，逼出了一个当时没人愿意接受的结论。
\end{zhengju}

今天我们知道，转座元件在自然界普遍得惊人。上一节算过账：人类基因组的将近一半是它们留下的。绝大多数拷贝早已突变失灵，成了动不了的“残骸”；但少数仍然活跃。比如 \emph{LINE-1} 元件，至今还会在人类的早期胚胎里跳到新位置——跳到无关紧要的区域，风平浪静；万一跳进一个关键基因，就可能引发疾病，某些血友病和癌症的病例，病根就是一次插入事件。

但演化对旧零件从不浪费。有些跳进来的序列，恰好带有一段能控制基因开关的短序列，被宿主“收编”成了调控元件——本来是自私的寄居者，最后变成了说明书里的标点符号。基因组里“自己搬家的段落”，就这样在“捣乱”和“被征用”之间，成了变异的重要来源。

\section{向邻居借一页：水平基因转移}

到目前为止讲的变化，都发生在“一家人”内部：从亲代传给子代，这叫垂直传递。但生物界还存在一条横向的通道：直接把基因\textbf{借}给没有亲缘关系的邻居。这叫\textbf{水平基因转移}。

细菌是借钱借得最凶的。它们能通过直接接触传递一小环 DNA（质粒），能从环境里捡起死细菌散落的 DNA 片段，还能靠噬菌体（感染细菌的病毒）在搬家时夹带私货——三种机制的细节在第 82 章讲基因工程时还会用到，这里只需记住：对细菌来说，基因组不是一家一户的私产，而是一个可以互相传阅的公共图书馆。

这件事跟你我的关系极其直接：抗生素耐药性。一个细菌偶然突变出分解抗生素的酶，它的后代当然都有耐药性——这是第 76 章讲过的垂直故事。但更麻烦的是，这个耐药基因如果恰好位于质粒上，这个细菌还能把它\textbf{当面转交}给另一个种类完全不同的细菌。于是医院里出现的“超级细菌”，往往不是靠慢慢突变熬出来的，而是几种耐药基因在细菌之间传来传去、层层打包的结果。这就是为什么医生反复强调：抗生素必须足量足疗程，能用窄谱就不用广谱——每一次滥用，都是在给这个“基因借阅网络”做选拔训练。

水平基因转移不只发生在细菌之间，偶尔也跨越大得惊人的鸿沟——这就是开场那只烤红薯的秘密。农杆菌天生就是“基因工程师”：它感染植物时，会把自己的一段 DNA（叫 T-DNA）插进植物细胞的基因组，让植物细胞为它制造养分。人类利用它这个本领做转基因（第 82 章），但 2015 年的研究发现，农杆菌\textbf{早就替自己做过了}：所有受检的栽培红薯品种里，都带着农杆菌的 T-DNA 片段，而野生近缘种里没有。这强烈暗示：正是这段“借来”的细菌基因，参与了红薯块根膨大这一关键性状的形成——人类几千年来选中并喜爱的红薯，也许正是大自然一次水平基因转移的产品。

需要说清楚的是：这种事在细菌里是家常便饭，在动植物身上则非常罕见。多细胞生物有把生殖细胞和体细胞分开的“围墙”（第 75 章），外来 DNA 就算钻进你脚趾的一个细胞，也进不了卵子和精子，传不下去。红薯的幸运，在于那次插入恰好发生在能开花结果、还能无性繁殖（插藤）的材料上。

\section{病毒留下的化石：内源病毒序列}

水平基因转移还有一位特殊的快递员：病毒。

有一类病毒叫逆转录病毒（艾滋病病毒 HIV 就属于这一类），它们的遗传物质是 RNA，入侵细胞后先反转录成 DNA，再把这段 DNA\textbf{缝进}宿主的染色体——这是它们繁殖策略的一部分，从此细胞每次复制 DNA，都会顺带帮病毒抄一份。

现在设想一个罕见但必然发生的意外：假如病毒缝进去的那个细胞，恰好是一颗生殖细胞，而这颗生殖细胞又恰好发育成了一个新个体——那么病毒的这页“广告单”，就会出现在这个新个体的每一个细胞里，并且传给他的所有后代。这样的残页叫\textbf{内源性逆转录病毒序列}。前面算过，它们约占人类基因组的 8\%。换句话说，你的基因组里躺着几千万年来一次次病毒感染的“案发现场”，层层叠叠，像地质岩层一样记录着祖先遭遇过的瘟疫。

这些残页绝大多数早已残破失能。但演化的修补匠又一次废物利用，而且这一次的成果堪称震撼：你之所以能坐在妈妈的子宫里出生，可能要感谢一段病毒基因。

胎盘的关键结构，是胎儿一侧的细胞会\textbf{融合}成一整层多核的“合胞体”，把母体和胎儿的组织隔开又连通，养分和氧气在这里交换。让细胞融合的那把“分子拉链”，是一个叫合胞素（syncytin）的蛋白质——而它的基因，正是一段古老逆转录病毒留下的残页：病毒原本用它来欺骗细胞膜、钻入细胞，哺乳动物把它收编过来，反过来用它建造胎盘。后续研究还发现，不同哺乳动物支系收编的合胞素来自\textbf{不同的}古老病毒，相当于同一件工具被独立地“捡”了好几次。一段致病病毒的遗物，就这样成了胎生动物演化的关键零件。

顺便说句公道话：病毒不只是借进，偶尔也会借出——把宿主的基因裹进自己的颗粒带给下一个宿主。在微生物的世界里，病毒可能是地球上规模最大的“基因搬家公司”。

\section{不动零件，改说明书：调控序列的变化}

前面几节讲的都是“序列本身”的变化。但还有一条更省事的进化捷径：零件一个不换，只改\textbf{使用说明}。

回忆第 74 章和第 75 章：一个基因在哪些细胞里开工、开多快、开到几岁，都由它旁边的调控序列决定。调控序列也是 DNA，也会突变。而调控序列的一个小改动，带来的形态变化可能大得惊人——因为它改变的不是零件的质量，而是零件的\textbf{部署方案}。

经典案例是三刺鱼。生活在海里的三刺鱼长着一对巨大的腹鳍骨刺，用来防捕食者；大约一万年前冰期结束后，许多三刺鱼迁居到淡水湖泊里，而很多湖泊种群的腹刺退化消失了。基因比对发现，不同湖泊的鱼是\textbf{各自独立}丢掉腹刺的，但原因常常指向同一个基因：\emph{Pitx1}。这个基因本身没坏——它的编码序列完好无损，坏的是它旁边一小段“在腹鳍部位开工”的调控开关。为什么进化选择改开关而不是改基因？因为 \emph{Pitx1} 在别处还有工作（比如脑和胸鳍），把基因整个敲掉会伤及无辜；而只撕掉“腹鳍使用说明”这一页，副作用最小。这就是调控进化的精妙之处：它允许\textbf{局部}改造。

人和黑猩猩的基因编码序列相似度约 99\%，这个著名数字的另一面是：我们两个物种的巨大差异，很大程度上不是零件清单的差异，而是部署方案的差异——同样的零件，在大脑皮层多开工几代细胞，在手部发育中晚关几天工，结局就天差地别。第 73 章讲过的乳糖耐受，也是一个调控区突变（让乳糖酶基因“成年后别下班”）在不同人群中被独立选中的例子。

把这条线索再推远一步，会撞上演化生物学最深的一个发现：指挥果蝇长眼睛的主控基因，和指挥老鼠、指挥人长眼睛的，是\textbf{同一个}基因的直系后代（在果蝇里叫 \emph{eyeless}，在脊椎动物里叫 \emph{Pax6}）。把老鼠的 \emph{Pax6} 装进果蝇，果蝇照样长出眼睛——只不过是果蝇式的眼睛。果蝇的复眼和你的眼睛在结构上没有半点相似之处，祖上也各自独立，可“在该长眼睛的地方启动整套建造程序”的那个总开关，却在五亿多年前就已经造好，之后被所有后代沿用至今。进化一旦造出好用的开关，就再也不肯重造。

\section{修补匠，而不是设计师}

现在可以把这一章的画面拼起来了。法国生物学家弗朗索瓦·雅各布（François Jacob，就是第 70 章操纵子模型的提出者之一）在 1977 年写过一篇著名文章，标题的大意是：进化是个\textbf{修补匠}（tinkerer）。

工程师设计一台机器，会先画图纸，再按图纸挑选最合适的零件，从零制造。修补匠不这样：他手边只有一堆旧货——影印出来的多余零件、捡来的外来零件、废掉的病毒残骸——他的工作方式是敲敲打打、东拼西凑，看到旧零件能凑合出新用途就用上。羽毛最初可能用于保温，后来被“改行”用于滑翔和飞行；眼睛晶状体里透明的晶状体蛋白，许多本来是做别的事的蛋白质（比如某些酶），被“征用”来透光；你的声带、下巴、手腕里的骨头，个个都能在鱼类的鳃弓和鳍条里找到前身（第 78 章的同源结构说的就是这件事）。

这解释了生物界一个著名的“不优雅”：为什么设计如此多瑕疵。人的视网膜装反了，感光细胞贴在血管和神经的\textbf{后面}，光线要穿过一层“电线”才到达传感器，还在视野里留下一个盲点；喉返神经从脑出发，绕到主动脉弓下面再返回喉咙，在长颈鹿身上这趟冤枉路长达四米。没有任何工程师会这样画图，但修补匠别无选择：他只能改造祖先传下来的旧图纸，每一步改动都得让中间产物还能活。演化没有推倒重来的选项。

这也意味着“改造旧零件”有个推论：今天每一种生物的基因组，都是一部\textbf{考古层堆积}。最底层是几十亿年前的古老核心（复制、翻译的机器，所有生命共享），上面压着一次次影印、一批批搬家元件、一页页病毒残片、一层层被改写的开关。读懂这本书，等于同时读到一部生命的编年史——第 81 章会讲科学家怎样真的把它读出来。

\begin{wuqu}
“人类基因组 98\% 是垃圾 DNA，都是没用的废物。”——这句话两头都错。先说“98\%”：它只说明这些序列不编码蛋白质，而“不编码蛋白质”远不等于“没用”——调控开关、各种功能性 RNA、染色体的结构件（着丝粒、端粒）都藏在这 98\% 里，本章讲的合胞素和被收编的转座元件开关也是。但反过来，也别矫枉过正地说“垃圾 DNA 全都有功能”：基因组里确实堆着大量真正惰性的残骸——失活的病毒残页、突变失灵的影印副本——它们大多只是被复制机器顺手照抄，既不帮忙也不碍事。真实的图景是：这 98\% 是一个杂物间，里面有工具、有备件、也有真正的垃圾，而演化正是靠翻这个杂物间找到新点子的。听到有人用“垃圾 DNA 其实都有用”来宣布某个惊人结论，或者用它来证明基因组“被完美设计”，你都要多问一句：证据是什么？
\end{wuqu}

\begin{shiyan}
用纸和笔模拟“影印、搬家和借用”，亲身感受变异怎样积累。\textbf{材料}：一张纸、一支笔、一个计时器（手机即可）、一段 30 字左右的句子（比如从这本书里抄一句），再找一位家人或同学帮忙。\textbf{步骤}：（1）把句子给对方看 10 秒钟后收起，请对方默写；再把你写的和原句对照，数出有几处“突变”。（2）玩“复印机游戏”：第一个人默写一句，把默写结果给第二个人照着默写，依次传 5 个人，最后和原句比较。（3）模拟“借用”：在第 3 个人的句子中间，硬塞进一句从别的书里抄来的话，再传给第 4 个人，看后面的人会不会把外来句当成原文的一部分继续抄。（4）模拟“影印”：把某个短语在原句里连写两遍，看它在后续传抄中保留还是消失。\textbf{观察点}：哪类变化最容易保留下来？被塞进的外来句在几轮传抄后“看起来”还像外来户吗？\textbf{想一想}：真实的 DNA 复制有校对和修复（第 72 章），传抄游戏没有——如果允许每人对照“原句的一个副本”改错，结果会怎样不同？\textbf{安全提示}：纯纸笔活动，无风险；别用太长的句子，30 字以内效果最好。
\end{shiyan}

\section{边界：这个比喻在哪里失效}

“基因组是一本书”这一章用得很顺手，但现在是拆它台的时候了。书是死物，基因组不是。至少有三处，比喻会把你带偏：第一，书不会自己决定哪页爱突变——真实的基因组里，不同区域的变动速率相差成百上千倍，修复系统的覆盖、序列本身的化学性质都在影响“哪里容易出事”，变动并不是均匀撒胡椒面。第二，书页之间的借用没有成本——而真实的水平转移和影印，绝大多数对携带者毫无好处甚至有害，能在群体中扩散开来的只是极少数；第 77 章讲过的遗传漂变还会告诉你，连“有没有用”都不是留存的唯一标准，运气也占很大一份。第三，也是最容易被忘掉的：本章讲的每一桩“基因组大事”，主角都是\textbf{群体}在世代间的积累，不是某个个体的一生。你个人的基因组这辈子基本定型；会影印、会借用、会改造的，是绵延亿万年的那条传代链条。

\section{回到那只烤红薯}

现在重新掰开那只烤红薯。它甜，是因为加热时淀粉被切成小段、又在酶的作用下变成糖（第 56 章的故事）；它面，是因为块根里堆满了淀粉颗粒——而这份“把养分囤进块根”的本领，很可能与亿万年前农杆菌硬塞进来的那几页外来基因有关。那几页基因没有走任何实验室流程，却完成了一次教科书级的水平基因转移：插入、保留、扩散到所有栽培品种，最后被人类的选择放大到每一块田里。

所以下次听到“转基因”三个字，你至少该想起三件事：第一，基因在物种之间流动，是大自然干了亿万年的老手艺，人类只是学会了按图纸施工（第 82 章）；第二，一段外来的、看似“非我族类”的序列，完全可能变成新物种性状的关键零件——你我的胎盘就是明证；第三，判断一项技术该问的问题，从来不是“天然还是人工”，而是“改了什么、证据如何、剂量多少”。

\section{还没讲完的故事}

这一章讲的是基因组“硬件”层面的翻修：影印、搬家、借页、改开关。但还有一笔最大的“借”，我们留到了下一章：大约二十亿年前，一个细胞把另一个细菌整个吞了下去，没有消化，反而让它住了下来——那就是你每个细胞里的线粒体。整个细胞器都是“借”来的，这算是水平基因转移的极限形态，还是一种全新的合作模式？第 80 章，我们进入合作、共生与多层次选择的世界。再往后，第 81 章讲人类怎样读出这些层层叠叠的古老文本，第 82 章讲我们怎样拿起农杆菌和细菌免疫系统留下的工具，开始自己动手改写这本书。

\begin{tiaozhan}
一段基因被影印成两份之后，命运会怎样？群体遗传学给出三条主路，可以粗粗估算它们的概率。设种群有效大小为 $N$，一个中性变异在群体中扩散到全体的概率约为它当前的频率——一个刚出现的新拷贝，频率是 $\frac{1}{2N}$，小到可以忽略。所以第一条路，也是绝大多数重复拷贝的下场：在几代人的时间里悄悄丢失，或者积累破坏性突变变成假基因。第二条路：两份拷贝把原来的工作量\textbf{平分}——比如原来一个基因要在两处开工，影印后正好一人管一处，这叫亚功能化，只要两份拷贝各自坏掉一个开关就可能被双双保留。第三条路最罕见也最激动人心：一份拷贝撞出全新的有用功能，被自然选择牢牢抓住（新功能化），红绿视蛋白和珠蛋白家族走的就是这条路。三条路的分流比例，至今仍是演化基因组学里活跃的研究问题；可以确定的是，“影印→新功能”绝不是必然剧本，而是百万次影印里偶尔中的一次彩票——但进化买彩票的本钱，是亿万年和亿万个体。
\end{tiaozhan}

\section*{练一练}

\begin{sikaoti}
\item 红视蛋白基因和绿视蛋白基因序列相似度约 98\%，且并排位于 X 染色体上。请画出它们“由一个祖先基因影印而来”的示意图（两条染色体、错位、交换），并标出：影印之后，哪一份拷贝“获得自由”？为什么原件不太容易积累大突变？
\item 麦克林托克的玉米斑块可以用“基因自己跳来跳去”解释。请你提出一个\textbf{不包含跳跃元件}的替代解释，再说说看：她的哪一条观察（斑块的可预测遗传、染色体断裂点的移动）最难被你的替代解释说明？
\item 一位病人感染了一种“超级细菌”，它同时对三种抗生素耐药。用本章的概念说明：这三种耐药基因更可能是一次次独立突变出来的，还是通过质粒“打包借阅”来的？这两种假设各会预测什么样的流行病学证据？
\item 合胞素基因来自古代病毒。有人说：“所以人类是病毒变的。”请指出这句话错在哪一层，并用一段话准确表述病毒序列进入基因组的真实含义。
\item 三刺鱼丢掉腹刺时，演化“选择”破坏 \emph{Pitx1} 的调控开关而不是编码序列。请解释为什么这条路线代价更小，并预测：如果某个湖泊里腹刺消失是因为编码序列整个坏掉，这些鱼身上可能出现什么别的异常？
\item 翻开你这份基因组的“杂物间”：如果未来有一种生物把人类基因组里的某段 \emph{Alu} 元件“收编”成了有用的开关，这段元件还算不算“垃圾”？请用“功能由谁定义、在什么时候定义”来组织你的回答。
\end{sikaoti}
